% ============================================================
% DRAFT NOTE (results)
% Contains: The confirmation result and its characterization. Order:
% headline earliness curve; comparator table per checkpoint; calibration;
% stratified breakdown (task, compressor, ratio); catastrophic vs graceful;
% what the forecaster looks at (ablation interpretation, deferred from
% Section 4); secondary endpoints.
% Written now: Outline only. No numbers exist yet.
% Pending: Everything, from results/recovered/quality-risk-v1/.
% Sources: agent REPORT.md and confirmation_result/report.json
% ============================================================
\section{Results}
\label{sec:results}

\subsection{Earliness: how soon the risk becomes readable}
\todo{Figure: AUROC and calibration vs checkpoint t, HERALD vs action-clock
vs prevalence, with prompt-cluster bands. One sentence on where the curve
crosses the usefulness floor. This is the paper's main figure.}
\begin{figure}[t]
\centering
\figplaceholder{AUROC (left) and calibration slope or ECE (right) as a
function of checkpoint t, for HERALD, action-clock, and prevalence, with 95\%
prompt-cluster bands}{results/recovered/quality-risk-v1/confirmation\_result/}
\caption{\todo{Earliness curve caption.}}
\label{fig:earliness}
\end{figure}

\subsection{Against the comparators}
\todo{Table: per checkpoint, log loss / AUROC / Brier skill for HERALD,
action-clock, action-only, prevalence, with simultaneous bounds on the
paired differences. State the gate outcomes plainly.}

\subsection{Calibration}
\todo{Reliability diagram at t = 32 (and 128); slope, intercept, ECE with
bounds. Whether the gate held.}

\subsection{Where it holds and where it degrades}
\todo{Stratified grid: task x compressor x ratio at t = 32; call out the
fragile strata (heavy ratios, specific compressors, constrained-format
tasks) and whether any subgroup was harmed relative to action-clock. Full
grid to appendix.}

\subsection{Catastrophic versus graceful damage}
\todo{Earliness split by catastrophe flag. Expectation to test: catastrophic
runs become readable earlier (loops have a signature); graceful damage is
harder. Report both curves.}

\subsection{What the forecaster is looking at}
\todo{Interpretation of Table~\ref{tab:ablation}: which feature groups carry
the signal; whether the mechanism hypothesis (uncertainty, instability,
churn among alternatives) is supported; whether divergence-forecast features
mattered. Keep interpretation tied to the ablation, not to feature
importances alone.}

\subsection{Secondary endpoints}
\todo{Major damage (theta frozen; report value), catastrophe risk, magnitude
given damage, all descriptive, opened after the primary decision.}
